\documentclass[12pt]{article}
\usepackage[utf8]{inputenc}
\usepackage[T1]{fontenc}
\usepackage{amsmath,amssymb}
\usepackage{geometry}
\geometry{margin=2.5cm}

\begin{document}

\section*{Lecture 3 -- Page 3}

\noindent\textbf{Definition.} Let $(S,\varphi)$, $\varphi : S \times S \to S$, $S \neq \varnothing$.

\begin{enumerate}
\item[(1)] If $\varphi$ is associative $\Rightarrow$ $(S,\varphi)$ is called a \textbf{semigroup}.

\item[(2)] A pair $(M,\varphi)$, $M \neq \varnothing$, $\varphi : M \times M \to M$.\\
If $\varphi$ is associative and has an identity element, $M$ is called a \textbf{monoid}.

\item[(3)] A pair $(G,\varphi)$, $G \neq \varnothing$, $\varphi : G \times G \to G$; if $\varphi$ is associative, has an identity element, and every element of $G$ is invertible with respect to $\varphi$ $\Rightarrow$ $G$ is called a \textbf{group}.
\end{enumerate}

\noindent\textbf{Examples:}
\begin{align*}
(\mathbb{N},+,0) &= \text{commutative monoid} \\
(M_2(\mathbb{R}), \cdot, I_2) &= \text{non-commutative monoid} \\
(\mathbb{Z},+,0),\ (\mathbb{R},+,0),\ (\mathbb{R}^*,\cdot,1) &= \text{commutative groups} \\
|A| \geq 2 \Rightarrow (T_A,\circ,1_A) &= \text{non-commutative monoid} \\
|A| \geq 3 \Rightarrow (S_A,\circ,1_A) &= \text{non-commutative group}
\end{align*}

\noindent\textbf{Computation rules in a semigroup/monoid/group}

\noindent\textbf{Def.} Let $S \neq \varnothing$, $(S,\cdot)$, ``$\cdot$'' associative, and $a \in S$, $n \in \mathbb{N}^*$.

\noindent We define
\[
a^n \overset{\text{def}}{=} \underbrace{a \cdot a \cdots a}_{n \text{ times}} =
\begin{cases}
a, & n=1 \\
a^{n-1}\cdot a, & n>1
\end{cases}
\]

\noindent In additive notation,
\[
na = \underbrace{a+a+\cdots+a}_{n \text{ times}} =
\begin{cases}
a, & n=1 \\
(n-1)a+a, & n>1
\end{cases}
\]

\noindent\textbf{Proposition.} Let $S \neq \varnothing$, $(S,\cdot)$, ``$\cdot$'' associative. Then:
\begin{enumerate}
\item[(1)] If $a \in S$, $m,n \in \mathbb{N}^*$ $\Rightarrow$ $a^m \cdot a^n = a^{m+n}$.
\item[(2)] If $a,b \in S$ and $ab=ba$ $\Rightarrow$ $(ab)^n = a^n \cdot b^n$, $(\forall)\, n \in \mathbb{N}^*$.
\end{enumerate}

\noindent\textbf{Proposition.} Let $(M,\cdot,e)$ be a monoid, $a \in M$, $a$ invertible.
\begin{enumerate}
\item[(1)] $a^n \cdot a^k = a^{n+k}$, $(\forall)\, n,k \in \mathbb{Z}$.
\item[(2)] $(a^n)^k = a^{n\cdot k}$, $(\forall)\, n,k \in \mathbb{Z}$.
\[
\left[\, a^0 \overset{\text{def}}{=} e, \qquad a^{-n} \overset{\text{def}}{=} (a^{-1})^n \,\right]
\]
\item[(3)] If $a\cdot b = b\cdot a$, $a$ and $b$ invertible $\Rightarrow$ $(ab)^k = a^k b^k$.
\end{enumerate}

\end{document}
